\section{Related Work}
\label{sec:related}

We organize prior work by the kind of evidence it provides for a
quantum-over-classical separation, and position our study against each group.

\para{Provable but non-simulable separations.}
\citet{liu2021rigorous} construct a binary classification task over $\mathbb{Z}_p^*$
whose labels are a half-interval in discrete-log space. Under the classical
hardness of the discrete log, no efficient classical learner exceeds
$\tfrac{1}{2}+1/\mathrm{poly}(n)$, while a quantum SVM with interval-state
feature maps provably reaches $\ge 0.99$ accuracy. \citet{jerbi2023shadows}
give an analogous RSA-style ``discrete cube root'' separation for classically
deployable shadow models, and \citet{glick2021covariant} show the discrete-log
task is a special case of covariant kernels. These results are the conceptual
bedrock, but they require fault-tolerant Shor execution on hundreds of qubits;
at simulable scale the number-theoretic problem is trivially solvable. Our
Experiment~B reproduces the \citeauthor{liu2021rigorous} task at $12$/$17$/$21$
bits and is explicit that it demonstrates \emph{off-the-shelf} classical failure,
not an asymptotic separation.

\para{Simulable empirical separations on engineered data.}
\citet{huang2021power} introduce the geometric difference \gdiff{} as a
pre-training diagnostic that upper-bounds achievable advantage, show that the
naive fidelity kernel concentrates in the exponentially large Hilbert space, and
propose the \emph{projected} quantum kernel built from reduced density matrices.
On engineered datasets that saturate \gdiff{}, the projected kernel beats every
tested classical model by $>20\%$ at up to $30$ qubits. \citet{jerbi2023beyond}
reach a similar conclusion for data re-uploading models, emphasizing that kernel
methods are not the right yardstick. Our Experiment~A reproduces the
\citeauthor{huang2021power} construction at $8$/$10$/$20$ qubits, and---unlike the
original, which reports a single advantage figure---we report a full
accuracy-vs-qubit-count curve, a within-quantum fidelity-vs-projected control, and
a paired significance test against the best non-kernel classical model.

\para{Hardware feasibility without a defeated baseline.}
\citet{havlicek2019supervised} introduce quantum kernel estimation and the
variational classifier on $2$ qubits; \citet{peters2021machine} run a fidelity
kernel on $10$--$17$ qubits of real hardware over $67$-dimensional supernova data;
and \citet{glick2021covariant} reach $27$ qubits with covariant kernels. These
demonstrate that quantum kernels are realizable, but quantum accuracy never
exceeds the classical baseline (and is often only compared to other quantum
models). We instead run noiseless $\ge 20$-qubit simulation specifically to
exhibit a separation, while disclosing that hardware shot noise may erode it.

\para{Cautionary and negative results.}
\citet{schuld2021models} prove every supervised quantum model is a kernel method,
so any separation must come from a classically hard kernel.
\citet{kubler2021inductive} show that sufficiently expressive quantum kernels have
exponentially small eigenvalues and cannot generalize unless they encode a
problem-specific bias---and that the favorable eigenvalues demand exponentially
many shots. \citet{schuld2021effect} characterize variational models as truncated
Fourier series, and \citet{tang2019quantum} dequantizes a flagship quantum ML
speedup. \citet{abbas2021power} argue for capacity advantages via effective
dimension, and \citet{du2020expressive,gao2022enhancing} give generative
expressivity separations that largely reduce to classically simulable tensor
networks. These results motivate our control experiment and our insistence on a
full classical suite: our fidelity-kernel collapse at 20 qubits is a direct
empirical instance of the concentration \citeauthor{kubler2021inductive} predict.
